% ---------------------------------------------------------------------------
% MSC PROPOSAL V1.2 — Cost-Aware Repository Change Localization with Sparse
% Impact Planning and Bounded Verification
%
% Historical version (V1.2). V1 is preserved immutably.
% Central story explicit; internal experiment codes removed from main prose;
% budget definitions frozen; factual corrections applied; bibliography
% verified against primary sources (arXiv API); no fabricated citations.
% ---------------------------------------------------------------------------
\documentclass[11pt,a4paper]{article}
\usepackage[T1]{fontenc}
\usepackage[utf8]{inputenc}
\usepackage{lmodern}
\usepackage{microtype}
\usepackage{amsmath,amssymb}
\usepackage{booktabs,tabularx,array}
\usepackage{enumitem}
\usepackage{graphicx}
\usepackage[margin=2.5cm]{geometry}
\usepackage{xcolor}
\usepackage{url}
\usepackage[colorlinks=true,linkcolor=blue,citecolor=blue,urlcolor=blue]{hyperref}
\emergencystretch=1.5em

\title{\textbf{Cost-Aware Repository Change Localization with Sparse Impact\\[2pt]
Planning and Bounded Verification}\\[6pt]
\large MSc Research Proposal (V1.2, supervisor-ready / print-candidate)}
\author{Ahmed Ehab}
\date{2026-09-16}

\begin{document}
\maketitle

\begin{abstract}
Large language models (LLMs) that edit a software repository must first decide
which files a requested change affects. This decision is costly: an explicit
file-by-file impact policy over a repository-scale candidate universe spends
hundreds of serialized records per call, and a second-stage reasoner that
revisits the scope can multiply inference cost by orders of magnitude. This
thesis asks a single central question:

\begin{quote}
\emph{Can a cheap/sparse first-pass repository localization policy followed by
bounded omission recovery achieve a better correctness--cost trade-off than
cheap localization alone and always-on expensive reasoning?}
\end{quote}

We investigate a representation that is cheap by construction ---
\emph{preserve-by-omission} sparse impact plans, which emit only non-preserve
decisions --- and evaluate it together with cheap lexical/static baselines and
an explicit \emph{bounded} second-stage omission-recovery mechanism, all under
matched inference budgets. Preliminary controlled and real-commit studies show
a large serialization-cost reduction for the sparse representation and a
meaningful zero-LLM lexical signal, but also that task-level omission risk is
not reliably predictable on current development evidence. The thesis therefore
evaluates the representation and the correctness--cost Pareto frontier, and
tests a pre-registered candidate-level bounded-verification mechanism, rather
than requiring a task-level risk scorer to succeed.
\end{abstract}

% ---------------------------------------------------------------------------
\section{Background, motivation, and problem}
\label{sec:background}

Repository-level coding agents and impact planners must select affected files
before editing. Two problems motivate the thesis. First, \emph{impact
inference}: did the model identify the right files? Second,
\emph{impact-policy representation}: can a complete file-action policy be
expressed without spending output budget on hundreds of repeated
\emph{preserve} decisions? Our controlled study shows that an explicit full
policy at a 16K completion cap serializes on average 144 decision records per
candidate universe, while a sparse policy that emits only non-preserve
decisions serializes on average 4.9 records (a $\sim$96.6\% reduction), with
large associated reductions in completion tokens and cost.

Second-stage reasoning is expensive. A graph-guided agent evaluated under a
shared protocol consumed on the order of a million tokens per task on the same
held-out tasks where our sparse first pass stays below a few thousand tokens.
If a cheap first pass already localizes the bulk of affected files, the
question is whether a \emph{bounded} second look can recover missed files
without paying the full always-on cost.

Given a repository revision $P$ (parent commit), a natural-language change
request $q$, and a candidate universe $U(P)$ of production files visible at
$P$, a localization system produces an affected-file set $\hat{S}(P,q)$. The
hidden observed change set $S^*(P)$ is used for evaluation only. The thesis
studies correctness--cost trade-offs under explicit budgets: representation
cost (sparse vs explicit) and bounded second-stage verification cost.

The observed historical diff is treated as an \emph{observed change-set proxy},
never as semantic ground truth (Section~\ref{sec:dataset}).

% ---------------------------------------------------------------------------
\section{Central thesis hypothesis and gap}
\label{sec:hypothesis}

\textbf{Hypothesis.} A sparse/cheap first-pass localization policy followed by
a bounded omission-recovery mechanism can dominate both cheap-localization-only
and always-on expensive reasoning on the correctness--cost Pareto frontier, at
least for the repositories and protocols we evaluate.

Everything else in the thesis supports this central story:

\begin{itemize}
  \item \emph{Preserve-by-Omission} = representation / first-pass contribution;
  \item \emph{BM25 / classical static change-impact analysis} = cheap baselines;
  \item \emph{graph-guided agent} = expensive system-level context;
  \item \emph{task-level omission-risk scorer} = investigated and currently
        rejected on development evidence;
  \item \emph{candidate-level bounded verification} = primary next mechanism;
  \item \emph{second and third repositories} = generalization.
\end{itemize}

Task-level RiskScorer success is \emph{not} a thesis requirement. Existing work
covers several individual components --- agentic graph-guided localization,
repository memory / retrieval-conditioned generation, explicit edit planning,
agentless localize-then-repair, and classical change-impact analysis. This
thesis evaluates \emph{their combination} as a sparse first pass plus bounded
omission recovery under matched budgets (see Section~\ref{sec:related}).

% ---------------------------------------------------------------------------
\section{Research questions and expected contributions}
\label{sec:rq}

\begin{description}
  \item[RQ1.] How does sparse explicit impact-policy representation affect
  inference/output cost and file-selection fidelity?
  \item[RQ2.] How accurately can cheap lexical/static methods localize impacted
  files relative to LLM planners under shared protocols?
  \item[RQ3.] Can bounded second-stage verification recover missed impacted
  files more efficiently than always-on or random matched-budget verification?
  \item[RQ4.] How do these correctness--cost trade-offs transfer across
  repositories and ecosystems?
\end{description}

If task-level risk modeling remains unsupported by evidence, RQ3 is answered by
the candidate-level mechanism; \emph{RiskScorer success is not required}.

\emph{Expected contributions:}
\begin{enumerate}
  \item Sparse explicit impact-plan representation with controlled evidence.
  \item A reproducible real-commit impact-localization dataset and protocol.
  \item Fair cheap/classical/LLM/agent comparison under common file-level
        metrics and explicit budgets.
  \item A bounded/selective-compute verification protocol (candidate-level by
        default; task-level only if evidence justifies it).
  \item Cross-repository correctness--cost evidence.
  \item Reproducibility artifacts and leakage-control methodology.
\end{enumerate}

No positive algorithmic result is promised.

% ---------------------------------------------------------------------------
\section{Positioning against related work}
\label{sec:related}

Existing work covers several individual components. We avoid any claim of the
form ``none combines\ldots'' unless a systematic matrix supports it. Our
positioning: this thesis evaluates the \emph{combination} of (a) an explicit
sparse impact-policy representation, (b) a cheap first pass, and (c) a
bounded omission-recovery mechanism, all under matched budgets, on real
commits across repositories.

Verified related lines (primary sources, 2026-09-16) and their relation:

\begin{itemize}
  \item \emph{Cost-aware LLM change-impact analysis:} requirements-level CIA
        under a cost lens (Etezadi et al., arXiv 2511.00262) --- same
        cost-aware question, requirements-to-requirements unit, not file
        localization.
  \item \emph{History / change-pattern augmentation of CIA:} Change-Patterns
        Mapping (TSE 2022), history-learned impact tools (ICSE-SEIP 2022),
        learning dependency-based predictors (IST 2015) --- the classical/
        history lineage that our Route~B history evidence must be positioned
        against.
  \item \emph{Modern transform + program-dependence-graph impact analysis:}
        Yan et al. (arXiv 2607.23355) --- graph+LLM impact analysis is active;
        our graph is one optional verifier.
  \item \emph{Selective-compute / efficiency+verification priors:} Repoformer
        (ICML 2024) and FastCoder (arXiv 2502.17139) --- selective retrieval and
        efficient verification at repository level (code-completion/generation
        tasks), guarding any generic ``selective escalation is novel'' claim.
  \item \emph{Graph-guided issue localization:} LocAgent, GraphLocator
        (arXiv 2512.22469), and LLM-driven iterative code-graph search
        (arXiv 2503.22424) --- the expensive graph-guided competitors.
  \item \emph{Recommendation-system CIA with industrial grounding:} Borg et al.
        (TSE 2017) --- classical lineage and industrial validity.
\end{itemize}

The comparison matrix (Section~\ref{sec:design}) records for each line whether
it has an explicit file policy, sparse omission representation, cheap first
pass, omission recovery, hard/matched budget, real-commit evaluation, and
cross-repository evidence. Where novelty is only defensible once experiments
support it, it is marked \emph{candidate novelty --- not yet claimed}.

% ---------------------------------------------------------------------------
\section{Proposed two-stage architecture}
\label{sec:architecture}

\begin{enumerate}
  \item \textbf{Sparse/cheap first-pass impact policy}: explicit file-level
        policy with preserve-by-omission serialization; cheap deterministic and
        LLM planner variants under the same contract.
  \item \textbf{Candidate-level bounded verification} (primary next mechanism):
        from the first-pass omitted candidates, rank suspicious candidates with
        independent structural / retrieval / history evidence, and verify only
        those under a hard per-task budget $B$. Compare with always-verify and
        random matched-budget verification.
\end{enumerate}

The flow:

\begin{center}
\begin{tabular}{@{}c@{}}
\hline
Change request \\
$\downarrow$ \\
Sparse / cheap first pass \\
$\downarrow$ \\
Omitted candidates \\
$\downarrow$ \\
Suspicion / ranking mechanism \\
$\downarrow$ \\
Bounded verification (budget $B$) \\
$\downarrow$ \\
Final affected-file set \\
\hline
\end{tabular}
\end{center}

% ---------------------------------------------------------------------------
\section{Dataset construction, labels, and leakage controls}
\label{sec:dataset}

\begin{itemize}
  \item \textbf{Real history, parent-only inputs.} Inference uses the parent
        commit $P$; the target $T$ is hidden. Live repository HEAD is never
        used (it moves and can leak future state).
  \item \textbf{Eligibility and exclusions.} Candidates must be single-parent
        commits with a meaningful natural-language intent, touching only
        production source files. Tests, configurations, documentation,
        migrations, and generated/vendor files are excluded from the candidate
        universe because they are not the target of file-level impact edits in
        this protocol. Leakage barrier: a case whose intent mentions a changed
        path is ineligible.
  \item \textbf{Deduplication.} Exact-proxy-set and shared-reference duplicates
        and suspected-related changes are removed with frozen deterministic
        rules, so accepted cases are independent.
  \item \textbf{Observed change-set proxy.} $P\!\to\!T$ changed production
        files are an \emph{observed} proxy, never semantic impact gold.
        Localization claims are scoped to historical-file recovery unless
        semantic adjudication (Section~\ref{sec:semantic}) supports more.
  \item \textbf{Splits.} A metadata-only, deterministic split (seed, hashes)
        is frozen before any model result. For djangoCMS V2-LARGE: development,
        an untouched internal test, and a reserve. The v1 40 cases are
        LEGACY-EXPOSED and never reused as untouched confirmation.
\end{itemize}

% ---------------------------------------------------------------------------
\section{Experimental design, baselines, budgets, metrics, statistics}
\label{sec:design}

\subsection{Experimental Design at a Glance}
\begin{table}[h]
\centering
\footnotesize
\resizebox{\textwidth}{!}{%
\begin{tabular}{@{}p{1.4cm}p{2.2cm}p{1.0cm}p{0.9cm}p{1.6cm}p{0.9cm}p{0.7cm}p{1.0cm}p{2.4cm}p{1.2cm}@{}}
\toprule
Study & Purpose & Repo & Split & Methods & Unit & Reps & Budget & Metrics & Status \\
\midrule
Representation & sparse-vs-full cost & curated & controlled & Full vs Sparse & task & 3 & cap 16K & tokens, records, cost, P/R/F1 & done \\
Real-commit & selection fidelity & djangoCMS & held-out & Full vs Sparse & task & 3 & cap 16K & P/R/F1/FNR & done \\
Graph agent & expensive context & djangoCMS & held-out & agent & task & 1 & per-call & P/R/F1, cost & done \\
Cheap baselines & zero-LLM signal & djangoCMS & dev & BM25/graph & task & -- & cheap & P/R/F1 & done \\
Task risk & risk predictability & djangoCMS & dev & features & task & 3 & -- & AUROC/AUPRC & done (neg.) \\
Candidate verify & omission recovery & djangoCMS & dev & rankers & candidate & -- & $B{=}5$ & FN-recovery & proposed \\
Cross-repo & transfer & Saleor/NestJS & TBD & frozen method & task & -- & frozen & per-repo & proposed \\
\bottomrule
\end{tabular}}
\caption{Experimental Design at a Glance (internal IDs appear in traceability notes only).}
\end{table}

\subsection{Baselines}
Random/cheap control, BM25, adaptive cheap retrieval (if justified),
classical/static change-impact analysis, Sparse planner, Full planner,
graph-guided reference agent (when reproducible), Sparse + Always Verify,
Sparse + Random Verify (matched budget), Sparse + Selective/Bounded Verify,
and an Oracle-routing upper bound. If task-level routing is unsupported, the
selective arm becomes candidate-level bounded verification.

\subsection{Related-work comparison matrix}
\centering
\footnotesize
\setlength{\tabcolsep}{3pt}
\begin{tabular}{@{}lccccccc@{}}
\toprule
Line & File pol. & Sparse rep. & Cheap 1st & Omission recov. & Matched budget & Real commit & X-repo \\
\midrule
Classical CIA (ripple/change-propagation) & no & no & yes & partial & no & some & no \\
IR localization & no & no & yes & no & no & some & no \\
RepoCoder / RepoAgent & no & no & no & no & no & some & no \\
CodePlan & yes & no & no & no & no & no & no \\
Agentless & no & no & partial & no & no & SWE-bench & no \\
LocAgent & no & no & no & no & no & yes* & no \\
GraphLocator & no & no & no & no & no & yes & no \\
Repoformer / FastCoder & no & no & partial & no & selective & no & no \\
This thesis & yes & yes & yes & yes & yes & yes & yes \\
\bottomrule
\end{tabular}

* LocAgent is evaluated here under a shared protocol (P5); not a faithful
reproduction of the published fine-tuned result.

\subsection{Frozen budget definition}
For zero-LLM candidate-ranking (Route B):
\begin{itemize}
  \item \textbf{Primary verification budget:} $B = $ number of omitted
        candidate files admitted to second-stage inspection.
  \item \textbf{Primary operating point:} $B = 5$ candidates per task.
  \item \textbf{Secondary:} $B \in \{1,3,10\}$.
  \item \textbf{Primary endpoint:} task-clustered FN recovery / missed-file
        recall gain at $B=5$.
  \item \textbf{Key control:} Sparse + Random Verify @ $B=5$.
  \item \textbf{Secondary:} revised FNR/F1, task coverage, scope inflation,
        runtime.
\end{itemize}
For a future LLM verifier: primary compute measure = additional total model
tokens per task; secondary operational measures = calls, USD cost, latency,
failures. Tokens, dollars, and latency are never collapsed into one metric.

\subsection{Multiple-comparison discipline}
One primary endpoint at $B{=}5$; one primary Selective-vs-Random comparison;
secondary endpoints clearly labeled; task-level paired bootstrap; no
feature-family winner chosen by dozens of post-hoc tests; exploratory analyses
reported as exploratory.

\subsection{Statistics}
Task is the independent unit; nested repetitions are never treated as
independent observations. Pre-registered class-balance gate for any risk
modeling: $\ge$10 independent positives and $\ge$10 negatives, a predeclared
feature above the random band, and direction stability across development /
validation. Bootstrap confidence intervals over tasks; random-risk bands at the
observed $n$; sample-size planning uses development evidence only.

% ---------------------------------------------------------------------------
\section{Preliminary evidence}
\label{sec:prelim}

All numbers are from audited studies in this repository; they are development
or controlled evidence and are reported with their bounds.

\begin{itemize}
  \item \textbf{Representation effect (controlled).} At a 16K completion cap,
        Sparse serialized on average \textbf{4.9} decision records vs
        \textbf{144.0} for Full (a $\sim$\textbf{96.6\%} reduction), with mean
        completion tokens $\sim$809 vs 8,383 and a large recorded-cost
        reduction. No universal semantic-superiority claim is made.
  \item \textbf{Historical real-commit study.} 60 cells over 10 held-out real
        djangoCMS commits (Full vs Sparse, 3 nested repetitions): selection is
        difficult, FNR is high, and there is no Full-vs-Sparse superiority
        claim (paired deltas' confidence intervals cross zero). Sparse
        execution is substantially cheaper.
  \item \textbf{Graph-guided agent comparison.} System-level shared-protocol
        context only: 402 calls / 32.8M tokens / $\sim$\$9.93 estimated on 10
        tasks, with 5/10 non-usable outcomes (2 true timeouts + 1
        context-length error + 2 completed-but-empty). Not a reproduction of
        the published fine-tuned result.
  \item \textbf{Cheap/static development baselines.} BM25 provides a
        meaningful zero-LLM localization signal on development data;
        graph-ranked seeds are comparable; K is an operating-point curve, not
        a fixed configuration. A classical change-impact (dependency
        propagation) baseline is implemented on 150 development cases.
  \item \textbf{Task-level omission-risk investigation.} Live Sparse
        development inference (90 cells; then 431 calls on 144 new development
        tasks with a fail-closed stop at the nominal 2.5M-token ceiling --- the
        final completed call produced a 1,964-token, $\sim$0.08\% post-call
        overshoot, and the run stopped before another call). A task-level
        RiskScorer is \emph{not} justified: the v1 candidate signal does not
        replicate within the new development cohort, and the apparent pooled
        signal is a candidate-universe-size artifact. No multivariable
        RiskScorer is fitted.
\end{itemize}

These results motivate a larger development dataset, protection of an untouched
internal test, and the bounded-verification design; they are not rewritten as a
success.

% ---------------------------------------------------------------------------
\section{Cross-repository generalization}
\label{sec:gen}

Stage 1 djangoCMS V2-LARGE (within-repo large evidence + untouched internal
test) $\to$ Stage 2 Saleor (different real system, same Python/Django stack)
$\to$ Stage 3 NestJS (cross-language / cross-framework evidence). Each stage
states its claim boundary: within-repo development + internal confirmation;
cross-repository confirmation; cross-ecosystem evidence. NestJS is preferred
but not forced; a predeclared eligible-pool yield criterion and a backup
protocol are frozen.

% ---------------------------------------------------------------------------
\section{Threats, contingency, and timeline}
\label{sec:threats}

\subsection{Threats to validity}
\begin{itemize}
  \item \textbf{Internal:} the observed proxy is not semantic gold; leakage
        control (parent-only inputs, intent-path-leakage off) is mandatory; no
        hidden-test peeking.
  \item \textbf{External:} $n$ and single-repository generalization; V2 and
        cross-repository evidence are the mitigation, with per-repository
        primary metrics.
  \item \textbf{Statistical:} low negatives in the development label;
        multiple comparisons; wide random bands --- addressed by
        pre-registration, class-balance gates, and bootstrap uncertainty.
  \item \textbf{Representation:} LLM serialization cost spans providers; we
        report both recorded and frozen-pricing cost and keep the provider
        pinned as far as allowed.
\end{itemize}

\subsection{Semantic-proxy validity (preparation)}
\label{sec:semantic}
Because the historical diff is only a proxy, we prepare a stratified
\emph{development-only} sample (20--30 tasks) with evidence packets (public
intent, parent context, $P\!\to\!T$ diff, changed production files, candidate
universe metadata) for later human adjudication. Human questions: was each
changed production file semantically necessary; is there evidence of
semantically impacted but unchanged files; is the commit mixed with incidental
refactor. Machine-assisted notes are prepared but are never called semantic
gold. Protocol: \texttt{docs/SEMANTIC\_PROXY\_AUDIT\_PROTOCOL.md}.

\subsection{Contingency}
If candidate-level bounded verification does not beat random at the primary
budget, the negative result is frozen and reported; no complex verifier is
invented.

\subsection{Timeline}
\begin{tabular}{ll}
\toprule
\textbf{Window} & \textbf{Activity} \\
\midrule
2026-09/10 & Proposal V1.2; V2-LARGE freeze + build \\
2026-10 & V2 development inference; cheap + classical baselines; candidate-level verification \\
2026-11 & Bounded-verification experiments; Saleor protocol + inference \\
2026-12 & NestJS cross-ecosystem check; Pareto analysis \\
2027-01/02 & Analysis, replication, threat revisions \\
2027-03/04 & Write-up; internal-test confirmation \\
2027-05 & Submission-ready thesis \\
\bottomrule
\end{tabular}

% ---------------------------------------------------------------------------
\section{References}
\label{sec:refs}

The authoritative bibliography is \texttt{references.bib} (V1.2). References
were verified against primary sources (arXiv API) on 2026-09-16. Preprints are
marked. No fabricated or placeholder citations remain.

\end{document}